\documentclass[11pt,a4paper]{article}

%% ─── Packages ────────────────────────────────────────────────────────────────
\usepackage[utf8]{inputenc}
\usepackage[T1]{fontenc}
\usepackage{geometry}
\usepackage{booktabs}
\usepackage{array}
\usepackage{graphicx}
\usepackage{amsmath}
\usepackage{url}
\usepackage{setspace}

\geometry{margin=1in}
\onehalfspacing

%% ─── Metadata ─────────────────────────────────────────────────────────────────
\title{\textbf{Protocol-Guided AI Smoking Cessation Counselor:\\
Comparing AI-Generated and Human-Curated Retrieval-Augmented Generation}}

\author{
  Adi [Last Name]\textsuperscript{1} \and
  Ibrahim [Last Name]\textsuperscript{1} \\[0.5em]
  \textsuperscript{1}[Institution / Department]
}

\date{\today}

%% ═══════════════════════════════════════════════════════════════════════════════
\begin{document}
\maketitle

%% ─── Abstract ─────────────────────────────────────────────────────────────────
\begin{abstract}
Retrieval-Augmented Generation (RAG) has emerged as a promising approach
for grounding large language model (LLM) outputs in domain-specific evidence,
reducing hallucination in high-stakes clinical applications.
This paper evaluates two RAG configurations for a protocol-guided AI smoking
cessation counselor: (1) an \emph{AI-generated} question-answer dataset
($n=5{,}936$ pairs) and (2) a \emph{human-curated} dataset ($n=4{,}847$ pairs).
We assess faithfulness, answer relevancy, context precision, and context recall
using the RAGAS 0.4.x evaluation framework on a 100-question clinician-reviewed
test set.
Human-curated RAG achieves a mean faithfulness of $0.657 \pm 0.334$,
outperforming AI-generated RAG ($0.531 \pm 0.320$) by 24\% relative improvement,
and reducing the high-risk hallucination rate from 43\% to 34\%.
Compared to a web-scraping RAG baseline (WebRAG), dataset RAG is
$2.41\times$ faster at cold-start and $1.44\times$ faster when web content is
pre-cached, demonstrating a clear latency advantage for real-time deployment.
These results underscore the clinical value of professional dataset curation
and the practical benefits of vector-indexed retrieval over live web scraping
in time-sensitive counseling contexts.
\end{abstract}

\textbf{Keywords:} Retrieval-Augmented Generation, smoking cessation,
RAGAS, faithfulness, hallucination, LLM evaluation, clinical AI

%% ─── 1. Introduction ──────────────────────────────────────────────────────────
\section{Introduction}

Smoking remains one of the leading causes of preventable death worldwide,
and evidence-based cessation counseling is critical for reducing tobacco-related
morbidity~\citep{who2023tobacco}.
AI-powered counseling assistants have shown promise in extending access to
behavioral support, yet their clinical deployment raises important questions
about factual accuracy and reliability of generated responses~\citep{thirunavukarasu2023}.

Retrieval-Augmented Generation (RAG) addresses hallucination by grounding
LLM outputs in a domain-specific knowledge base~\citep{lewis2020rag}.
However, the quality of that knowledge base—whether compiled by AI or by
human experts—fundamentally determines the system's reliability.
Prior work by Ibrahim et al.\ established a WebRAG baseline achieving a mean
faithfulness of $0.86$ by retrieving live content from authoritative web sources.
Our study extends this by asking: \emph{does the origin and curation quality of
the underlying dataset matter, and at what latency cost?}

We compare three configurations evaluated against the Gemini Protocol—a
structured smoking cessation framework developed by Dr.\ Louis Willis:
\begin{enumerate}
  \item \textbf{Baseline LLM}: GPT-4o-mini without any retrieved context.
  \item \textbf{AI-Generated RAG}: Vector retrieval from a 5,936-pair
        AI-generated Q\&A dataset indexed in ChromaDB.
  \item \textbf{Human-Curated RAG}: Vector retrieval from a 4,847-pair
        professionally curated Q\&A dataset.
\end{enumerate}

%% ─── 2. Related Work ──────────────────────────────────────────────────────────
\section{Related Work}

\paragraph{Hallucination in clinical LLMs.}
LLMs are prone to generating plausible-sounding but factually incorrect
statements, a phenomenon termed hallucination~\citep{ji2023survey}.
In clinical settings, hallucination can mislead patients and undermine trust,
making faithfulness evaluation essential~\citep{singhal2023}.

\paragraph{Retrieval-Augmented Generation.}
RAG~\citep{lewis2020rag} conditions LLM generation on retrieved documents,
substantially reducing hallucination rates compared to generation-only baselines.
Dense retrieval using vector databases such as ChromaDB enables efficient
semantic matching at inference time.

\paragraph{RAGAS evaluation framework.}
RAGAS~\citep{es2023ragas} provides reference-free and reference-based metrics
for evaluating RAG pipelines: faithfulness (are claims supported by context?),
answer relevancy (does the answer address the question?), context precision
(is retrieved context relevant?), and context recall (are reference answers
covered by retrieved context?).

\paragraph{Smoking cessation AI.}
Prior work has demonstrated that LLM-based counseling assistants can
deliver evidence-consistent advice~\citep{eapen2023},
with RAG improving adherence to clinical guidelines.

%% ─── 3. Methods ───────────────────────────────────────────────────────────────
\section{Methods}

\subsection{Test Set}

We used the LFV-reviewed subset of Dr.\ Willis's 150-question smoking cessation
test set, excluding 23 questions flagged by a clinical reviewer, yielding
$n = 127$ clean questions. We report results on the first $n = 100$ questions.

\subsection{RAG Datasets}

\textbf{AI-Generated dataset}: 5,936 question-answer pairs generated by
GPT-4o-mini from the Gemini Protocol documentation.

\textbf{Human-Curated dataset}: 4,847 question-answer pairs professionally
curated and verified against the Gemini Protocol.

Both datasets were indexed in an ephemeral ChromaDB collection using
cosine similarity. Retrieval used top-$k = 3$ nearest neighbours.

\subsection{WebRAG Baseline}

Following Ibrahim et al., we also measure a web-scraping approach that
fetches content from five authoritative sources at query time:
\texttt{smokefree.gov}, CDC, Mayo Clinic, American Cancer Society,
and the American Lung Association.

\subsection{Answer Generation}

All configurations used GPT-4o-mini (temperature $= 0.3$, max\_tokens $= 300$)
with a smoking cessation counselor system prompt.
RAG configurations prepended retrieved context to the system prompt.

\subsection{Evaluation Metrics}

We report four RAGAS 0.4.x metrics:
\begin{itemize}
  \item \textbf{Faithfulness} ($F$): fraction of claims in the response
        that are entailed by the retrieved context $\in [0, 1]$.
        Hallucination rate $= 1 - F$.
  \item \textbf{Answer Relevancy}: semantic similarity between the response
        and a set of hypothetical questions generated from it.
  \item \textbf{Context Precision}: proportion of retrieved context chunks
        that are relevant to the ground-truth answer.
  \item \textbf{Context Recall}: coverage of the ground-truth answer by
        retrieved context.
\end{itemize}

\subsection{Latency Measurement}

Per-question wall-clock latency was measured using \texttt{time.perf\_counter()}.
Three conditions were timed: Dataset RAG (retrieval + generation),
WebRAG warm (cached web content + generation), and WebRAG cold
(fresh URL fetch per question + generation).

%% ─── 4. Results ───────────────────────────────────────────────────────────────
\section{Results}

\subsection{Faithfulness and Hallucination}

Table~\ref{tab:faithfulness} presents faithfulness statistics for all
configurations alongside Ibrahim et al.'s reference results.

\begin{table}[ht]
\centering
\caption{Faithfulness Comparison. Mean faithfulness $\in [0,1]$,
higher is better; hallucination rate $= 1 - \text{Mean}$;
\% High-Risk = proportion with faithfulness $< 0.5$.
}
\label{tab:faithfulness}
\resizebox{\textwidth}{!}{%
\begin{tabular}{lccccccr}
\toprule
\textbf{Configuration} & \textbf{Mean} & \textbf{SD} & \textbf{Median}
  & \textbf{95\% CI} & \textbf{Halluc.\ Rate} & \textbf{\% High-Risk} & \textbf{n} \\
\midrule
Ibrahim et al.: Base LLM & 0.58 & 0.22 & 0.61 & {[}0.42--0.74{]} & 0.40 & 20\% & 10 \\
Ibrahim et al.: WebRAG   & \textbf{0.86} & 0.11 & 0.88 & {[}0.84--0.88{]} & 0.14 & 4\% & 100 \\
\midrule
Ours: AI-Generated RAG   & 0.53 & 0.32 & 0.58 & {[}0.47--0.59{]} & 0.47 & 43\% & 100 \\
Ours: Human-Curated RAG  & \textbf{0.66} & 0.33 & 0.70 & {[}0.59--0.72{]} & 0.34 & 34\% & 100 \\
\bottomrule
\end{tabular}%
}
\end{table}

Human-curated RAG ($\bar{F} = 0.657$) outperforms AI-generated RAG
($\bar{F} = 0.531$), a difference of $\Delta = 0.126$ ($24\%$ relative improvement).
Both systems exhibit substantially higher variance than Ibrahim et al.'s
WebRAG ($\sigma = 0.33$ vs.\ $0.11$), reflecting the heterogeneous nature of
the static dataset compared to targeted live retrieval.

\subsection{Full RAGAS Metrics}

Table~\ref{tab:ragas} presents all four RAGAS metrics.

\begin{table}[ht]
\centering
\caption{Full RAGAS Metrics (mean $\pm$ SD). N/A: metric not applicable
or evaluation not completed.}
\label{tab:ragas}
\begin{tabular}{lcccc}
\toprule
\textbf{Configuration} & \textbf{Faithfulness} & \textbf{Answer Relevancy}
  & \textbf{Context Precision} & \textbf{Context Recall} \\
\midrule
Baseline LLM        & N/A            & $0.871 \pm 0.187$ & N/A              & N/A \\
AI-Generated RAG    & $0.531 \pm 0.320$ & $0.781 \pm 0.338$ & $0.463 \pm 0.476$ & $0.262 \pm 0.330$ \\
Human-Curated RAG   & $0.657 \pm 0.334$ & $0.822 \pm 0.289$ & N/A              & $0.395 \pm 0.368$ \\
\bottomrule
\end{tabular}
\end{table}

The baseline LLM achieves the highest answer relevancy ($0.871$),
likely because its responses are unconstrained by retrieved context and
optimised for question-answering fluency.
Human-curated RAG shows higher answer relevancy ($0.822$) than
AI-generated RAG ($0.781$), suggesting that higher-quality context
also improves response quality.

\subsection{Latency Comparison}

Table~\ref{tab:latency} reports per-question wall-clock latency across conditions.

\begin{table}[ht]
\centering
\caption{Latency Comparison (seconds per question, $n=100$).
Dataset RAG uses human-curated data.}
\label{tab:latency}
\begin{tabular}{lrrrrrr}
\toprule
\textbf{Method} & \textbf{Mean (s)} & \textbf{Median (s)} & \textbf{SD}
  & \textbf{p95 (s)} & \textbf{n} \\
\midrule
WebRAG (cold, per-Q scrape) & 7.697 & 7.595 & 2.245 & 12.362 & 100 \\
WebRAG (warm, cached)       & 4.620 & 4.275 & 2.571 &  9.572 & 100 \\
Dataset RAG (human-curated) & \textbf{3.198} & \textbf{2.866} & 1.605 & 6.316 & 100 \\
\bottomrule
\end{tabular}
\end{table}

Dataset RAG is $2.41\times$ faster than cold WebRAG and $1.44\times$ faster
than warm (pre-cached) WebRAG. This advantage stems from ChromaDB's
sub-second vector retrieval ($\bar{t}_\text{retrieval} \approx 0.15$\,s),
compared to URL fetching which averages $\approx 3$--5\,s even with cached content.

%% ─── 5. Discussion ────────────────────────────────────────────────────────────
\section{Discussion}

\paragraph{Faithfulness gap relative to WebRAG.}
Both of our RAG systems achieve lower faithfulness than Ibrahim et al.'s
WebRAG ($0.86$).
This gap reflects a fundamental architectural difference: WebRAG retrieves
targeted, query-specific content from authoritative live sources, while our
systems retrieve from a fixed indexed corpus.
We expect that supplementing our dataset with periodic content updates
from authoritative sources, or adopting a hybrid retrieval strategy
(dataset + WebRAG fallback), would narrow this gap substantially.

\paragraph{Value of human curation.}
The 24\% relative improvement in faithfulness from AI-generated to human-curated
RAG, and the 9-percentage-point reduction in high-risk hallucination rate
(43\%~$\to$~34\%), quantify the clinical value of professional dataset curation.
Given the sensitivity of smoking cessation counseling—where inaccurate advice
may undermine quit attempts or patient trust—this quality improvement
justifies the curation cost.

\paragraph{Latency implications.}
Dataset RAG's latency advantage ($2.41\times$ vs.\ cold WebRAG) is
practically significant for real-time chatbot applications where
response times above 5--10\,s degrade user experience.
Even in a warm-cache scenario, Dataset RAG maintains a $1.44\times$ advantage,
making it preferable for high-throughput deployment without requiring
per-session web scraping infrastructure.

\paragraph{Limitations.}
(1) Context precision for human-curated RAG was not obtained due to
OpenAI API rate limits during evaluation; a rerun is scheduled.
(2) The test set ($n=100$) is drawn from a single clinical protocol,
limiting generalisability to other smoking cessation frameworks.
(3) GPT-4o-mini was used for both answer generation and RAGAS scoring,
which may introduce model-specific biases.

%% ─── 6. Conclusion ────────────────────────────────────────────────────────────
\section{Conclusion}

We present a comprehensive evaluation of AI-generated versus human-curated RAG
for a protocol-guided smoking cessation counseling system.
Human curation provides a meaningful improvement in faithfulness (0.53~$\to$~0.66)
and reduces high-risk hallucination rates (43\%~$\to$~34\%) compared to
AI-generated datasets.
Both systems are substantially faster than live WebRAG approaches,
with Dataset RAG achieving $2.41\times$ the speed of cold WebRAG and
$1.44\times$ the speed of warm WebRAG.
Future work will explore hybrid retrieval (dataset + live web sources),
expand the evaluation corpus beyond a single protocol, and investigate
fine-tuned embeddings for domain-specific semantic matching.

%% ─── Acknowledgements ─────────────────────────────────────────────────────────
\section*{Acknowledgements}
The authors thank Dr.\ Louis Willis for providing the clinician-reviewed
test set and the Gemini Protocol documentation.

%% ─── References ───────────────────────────────────────────────────────────────
\bibliographystyle{plainnat}

%% Replace with your .bib file or expand inline references:
\begin{thebibliography}{9}

\bibitem{who2023tobacco}
World Health Organization (2023).
\textit{Tobacco Fact Sheet}.
\url{https://www.who.int/news-room/fact-sheets/detail/tobacco}

\bibitem{thirunavukarasu2023}
Thirunavukarasu, A.J., et al.\ (2023).
Large language models in medicine.
\textit{Nature Medicine}, 29, 1930--1940.

\bibitem{lewis2020rag}
Lewis, P., et al.\ (2020).
Retrieval-augmented generation for knowledge-intensive NLP tasks.
\textit{Advances in Neural Information Processing Systems}, 33, 9459--9474.

\bibitem{es2023ragas}
Es, S., et al.\ (2023).
RAGAS: Automated evaluation of retrieval augmented generation.
\textit{arXiv preprint arXiv:2309.15217}.

\bibitem{ji2023survey}
Ji, Z., et al.\ (2023).
Survey of hallucination in natural language generation.
\textit{ACM Computing Surveys}, 55(12), 1--38.

\bibitem{singhal2023}
Singhal, K., et al.\ (2023).
Large language models encode clinical knowledge.
\textit{Nature}, 620, 172--180.

\bibitem{eapen2023}
Eapen, B.R., et al.\ (2023).
Proof of concept for using multimedia AI for patient education: ChatGPT and
YouTube.
\textit{Cureus}, 15(2), e35511.

\end{thebibliography}

\end{document}
